% Français 9115, batch 20 — Gallica views 381–400.
% Pass: Opus 5.5 (claude-opus-5-5), 2026-09-22 — first pass, unchecked against the views by a human.
%
% What the twenty views hold. Twelve views carry writing, eight do not.
% Two pieces, and between them two leaves of figures and two small leaves
% of calculation; nothing here is sorted.
%
% (1) Views 381–389: the end of a fair copy (lightly corrected) of her
% French remarks on the vibrating lamina, the piece that follows Euler's
% six cases. Each leaf is written on one side only and photographed twice:
% 381, 383, 385, 387, 389 are the written sides; 380, 382, 384, 386, 388,
% 390 are the blank backs showing the writing through, reversed (382
% flopped was laid on 381 and matches line for line). View 381 opens in
% mid-sentence (« est fixée et l'autre appuyée, des valeurs relatives… »):
% the sentence begins on view 379, in batch 19. Seen through view 380,
% flopped, that leaf carries « (30) », « 185 » and a paragraph numbered 24
% (« J'ai déjà remarqué que, dans le cas où il y a un nombre impair de
% nœuds… »), and a béquet at its head with « (29) » — seen only by
% transparency, so offered to the coordinator, not transcribed.
%   v. 381  the six cases of Euler reduce to four independent ones (I, III,
%           IV, VI, §§ 15, 30, 35, 43); case II (§ 23) agrees with case I
%           with ω for ½ω, the half-lamina with one end supported being half
%           of a lamina free at both ends; case V (§ 39) and case VI (§ 43)
%           should stand in the same relation and seem not to.
%   v. 383  the difficulty cleared: Euler takes the origin at the supported
%           end in case V; moving it (1−x for x) turns the ordinate of case
%           VI into the expected form. Then the check by Euler's conditions
%           for a simply supported end (y = 0, ddy/ds² = 0, § 13).
%   v. 385  why the two halves of a lamina free at both ends, with no node
%           at the middle, cannot be treated as whole laminae (her § 14);
%           § 25: the square plates of Chladni, free on every side, which in
%           several cases can be treated as four plates under the same
%           conditions (Chladni's Traité d'acoustique, § 110); a marginal
%           note of hers keyed by a cross.
%   v. 387  Chladni's figures 63, 68 a, 75 and 82 as one figure repeated
%           4, 9 and 16 times; the plate of side am divided into mm squares
%           of side a; the symmetry of the sections through the points α,
%           α + a, … and α', α' + a, ….
%   v. 389  α = α' = 0 for stretched membranes, by Biot's memoir (§ 24,
%           tome 4 of the Mémoires de l'Institut); fifteen cases for the
%           square plate; a reflexion on the word « appuyé ». The leaf ends
%           mid-sentence (« que dans la »), squeezed against the right edge,
%           and no view of this batch continues it.
%
% (2) View 391: a leaf of figures, ink 191, no text: two laminae EF with
%     their nodes (« fig 1 », « fig. 1′ », the figure read 1 or 7), and
%     Chladni's figures 63, 68 a, 75, 82 redrawn, as cited at view 387.
%     Views 392 and 394: two small leaves, ink 192 and 193, of
%     calculation in e^{xω}, sin xω, cos xω with double signs —
%     « resultat absurde », then the condition dy/dx = 0 at x = ½ — with
%     no heading and no link to (1). View 393 is blank: through it show,
%     reversed, the calculations of 392 at the top and figures 75 and 82
%     of 391 at the foot; view 395 is the
%     blank back of 394 (its calculations and the library stamp show
%     through). View 396: a second leaf of figures, ink 194, six square
%     plates (« fig. 67 c », « fig 72 a », « fig 78 », « fig. 85 », and two
%     unnumbered), none cited by a text of this batch; view 397 is its
%     blank back.
%
% (3) Views 398–400: the opening of a new piece, a heavily worked draft of
%     her prize memoir, « Mémoire sur la question proposée par la premiere
%     classe de l'institut “Donner la théorie mathématique des vibrations
%     des surfaces élastiques et la comparer à l'experience” », with an
%     epigraph from Bacon, Novum Organum I, XCII, in Latin. Page (1), view
%     398, is the preamble (the plane surface, the regular motion like a
%     simple pendulum, Chladni's experiments, the supposition z very small
%     carried to curved surfaces); page (2), view 399, is the back of the
%     same leaf — § 1 « Recherche de l'équation différentielle dela surface
%     élastique vibrante », N.o 1, the sum of the moments ∫∫Eβ de, the angle
%     e between two planes; page (3), view 400, on the next leaf, the
%     coordinates of four neighbouring points and the values of A, B, C.
%     It runs on to page (4), view 401, in batch 21.
%
% Views not transcribed, all blank backs showing the recto through,
% reversed: 382, 384, 386, 388, 390, 393, 395, 397. No repeated
% photograph, no microfilm target.
%
% The numbers on the leaves, and why no \folio{} is written.
% (a) Top right, in ink, a broad pen, one per recto, the running series of
% the earlier batches: 186 (v. 381), 187 (383), 188 (385), 189 (387), 190
% (389), 191 (391), 192 (392), 193 (394), 194 (396, last digit against the
% edge), 195 (398), 196 (400). None on the blank backs and none on 399,
% which is the written back of 195. Unbroken, eleven leaves. As in every
% batch of this volume, the figures are penned with thick and thin
% strokes, and no \folio{} is written (hand.md, « The open question of this
% volume »). (b) Top centre, in parentheses and struck through with a
% horizontal stroke, her own page numbers of piece (1): (31) v. 381, the
% second digit doubtful, (32) 383, (33) 385, (34) 387, (35) 389. They are
% not the underlined, unstruck 30–38 of views 283–299 (batch 15): another
% series. (c) Her paragraph numbers in piece (1): 25 at v. 385, written in
% the line; a cross-reference to her « § 14 » at v. 385. (d) The pages of
% piece (3): (1) struck at v. 398, (2) at 399, (3) at 400, not struck.
% (e) Chladni's figure numbers (63, 68 a, 75, 82; 67 c, 72 a, 78, 85) and
% Euler's section numbers (§§ 13, 15, 23, 30, 35, 39, 43) cited in the
% text. (f) On the large mounting sheets behind the small leaves (v. 394,
% 395, 397), two short oblique grey strokes, not a figure. No cited folio
% of the literature was found; nothing here is Laubenbacher and
% Pengelley's « Manuscript C ».
%
% Notation, kept: ω the angle of Euler's cases, x the fraction of the
% length, α β γ δ his coefficients (her α looks like « ∝ », her γ like
% « v »); e the exterior angle between two planes, β the small line of
% intersection, E the force, underlined letters as she underlines them
% (\underline{}); A, B, C and A', B', C' the coefficients of two planes;
% x', y', z' … x''', y''', z''' the coordinates of three points; d²z
% written with the 2 above the z. « sin », « cos », « tang. » set as
% \text{sin} etc., stops not recorded one by one. The dotted leaders on
% view 400 are hers and set as \ldots in maths. Her spelling kept:
% « équivallent », « parrallele », « précedant », « egals », « coordonnés »,
% « moitiées », « complette », « quarrées », « fesait », « pourrait » beside
% « pourroit », « transversales » (written so that it reads « transvertales »),
% dela, laquestion, cequi, aumilieu, audessus, cepoint, adire. Roman
% numerals of the cases as written, including the two apparently swapped
% at v. 381 and « cas VI » where V is meant; noted where they stand.
%
% What could not be read. \ill{}: the end of the last word of Bacon's
% reference (v. 398, it begins « lond »); in the heavily struck lines of
% v. 398, the end of « en substituant … », the middle of the line
% « différentielles … les différentielles », the rest of « de la courbe
% … », the end of « exprime-roient … », and a blotted word between
% « qu'elle » and « appréciables » in the margin. Doubted in notes only,
% because they stand inside formulas: the second digit of « (31) », the
% « 2 » under the radical (v. 381), the blotted « e^ω) cos » closing
% v. 383, a blotted exponent (v. 394), the accents of (BC' − B'C) (v. 400).
% \uncertain{}: « 1 » under a blot and « 13 » (v. 383); « 14 »,
% « transversales », « supposer », « 110 », « contigue » (v. 385); « 63 »
% twice (v. 387); « 63 », « quinze », « transversales » (v. 389); the
% figure « 1 » twice and « 63 » (v. 391); « &c » (v. 392);
% « n'a fait qu'un », « lond » (v. 398). Nothing was guessed to fill a gap.
%
% Nothing here was checked against any published edition or study of
% Germain, Euler, Chladni, Biot or Bacon.
\documentclass[11pt,a4paper]{article}
% The preamble shared by every transcription of the mathematicians' manuscripts in Gallica.
%
% It rests on one idea: **uncertainty compounds, it does not resolve.** A
% transcription that smooths over a doubtful word has destroyed the only thing
% it was made for — a reader must be able to tell, at every point, what was on
% the page and what was supplied. The apparatus macros below are therefore the
% critical apparatus, not decoration.
%
% The same commands are recognised by `scripts/render.mjs`, which produces the
% site's reading view, and by `scripts/tei.mjs`. A macro added here must be
% added there too, or rendering fails loudly — which is the wanted behaviour.
%
% Comments here are English, like the rest of the repository. The documents
% this preamble sets are French, and that is deliberate: see the skills.

\ProvidesPackage{germain}[2026/09/15 Transcriptions of mathematicians' manuscripts in Gallica]

\usepackage{amsmath,amssymb,amsthm}
\usepackage{mathrsfs}
% Kept from the parent preamble so the renderer's diagram support stays
% available; these manuscripts draw no commutative diagrams, and nothing here uses it.
\usepackage{tikz-cd}
\usepackage[english,french]{babel}
\usepackage{fontspec}

% --- Greek ------------------------------------------------------------------
% The text face stays fontspec's default, Latin Modern, which has no Greek:
% XeTeX drops every glyph it lacks with a line in the log and nothing on the
% page, and the Greek words the manuscripts quote vanished from the PDFs. So
% the Greek and Greek Extended blocks get a XeTeX character class of their own,
% and entering or leaving a run of it switches the family to CMU Serif — the
% same Computer Modern design, with polytonic Greek — and back. Only the
% family changes: a Greek word in \textit or \textbf keeps its shape and series.
% The transitions to and from every other class carry the tokens that class
% has towards an ordinary letter, so French spacing before « ; » or inside
% « » still holds after a Greek word. No grouping, so no brace or \endgroup
% can come out unbalanced; maths is left alone, since XeTeX inserts nothing
% there. Kept identical in `exercices.sty`.
\makeatletter
\newfontfamily\germainGreekfont{cmun}[
  Extension=.otf, NFSSFamily=germaingreek,
  UprightFont=*rm, ItalicFont=*ti, SlantedFont=*sl,
  BoldFont=*bx, BoldItalicFont=*bi, BoldSlantedFont=*bl]
\newXeTeXintercharclass\germain@greek
\def\germain@greekrange#1#2{%
  \@tempcnta=#1\relax
  \loop
    \XeTeXcharclass\@tempcnta=\germain@greek
    \ifnum\@tempcnta<#2\advance\@tempcnta\@ne\repeat}
\germain@greekrange{"0370}{"03FF}
\germain@greekrange{"1F00}{"1FFF}
\def\germain@greekprev{\rmdefault}
\def\germain@greekin{%
  \edef\germain@greekprev{\f@family}\fontfamily{germaingreek}\selectfont}
\def\germain@greekout{\fontfamily{\germain@greekprev}\selectfont}
\def\germain@greekjoin{%
  \XeTeXinterchartokenstate=\@ne
  \@tempcnta=\z@
  \loop
    \ifnum\@tempcnta=\germain@greek\else
      \XeTeXinterchartoks\germain@greek\@tempcnta=\expandafter{%
        \expandafter\germain@greekout\the\XeTeXinterchartoks\z@\@tempcnta}%
      \XeTeXinterchartoks\@tempcnta\germain@greek=\expandafter{%
        \the\XeTeXinterchartoks\@tempcnta\z@\germain@greekin}%
    \fi
    \ifnum\@tempcnta<4095\advance\@tempcnta\@ne\repeat}
% babel-french sets its own classes, and saves and restores the interchar
% state, each time a language is selected: join again after it does.
\addto\extrasfrench{\germain@greekjoin}
\addto\extrasenglish{\germain@greekjoin}
\AtBeginDocument{\germain@greekjoin}
\makeatother

\usepackage{geometry}
\geometry{a4paper,margin=25mm,marginparwidth=28mm}
\usepackage{marginnote}
\usepackage{xcolor}
\usepackage[normalem]{ulem}
\setlength{\emergencystretch}{3em}
\usepackage{hyperref}
\hypersetup{colorlinks=true,linkcolor=black,urlcolor=black,pdfborder={0 0 0}}

% --- Batch metadata ---------------------------------------------------------
% Taken as they stand by the rendering script, which shows them at the head of
% the reading view. They fix what the file claims to cover. The macro names are
% the parent project's, so that the scripts stay shared; \folder here means the
% volume's slug — `fr-9115` — and \pages counts Gallica views.

\newcommand{\folder}[1]{\gdef\grothFolder{#1}}
\newcommand{\batch}[1]{\gdef\grothBatch{#1}}
\newcommand{\pages}[2]{\gdef\grothFirst{#1}\gdef\grothLast{#2}}
\newcommand{\foldertitle}[1]{\gdef\grothFoldertitle{#1}}

% The holder's own shelfmark, printed as they write it: « Français 9115 ».
\newcommand{\shelfmark}[1]{\gdef\grothShelfmark{#1}}

% The dating, copied verbatim from the holder's catalogue — never inferred
% here. « XIXe siècle » is the BnF's; a pass that dates a leaf from its content
% says so in a \note{}, as its own inference.
\newcommand{\dating}[1]{\gdef\grothDating{#1}}

% The status notice, printed in the title block and repeated as a diagonal
% watermark on every page of the PDF. The manuscripts are in the public
% domain; what this line declares is not a right but a quality — an
% unverified first pass by a machine. The skills write the line; a file
% without it gets no watermark, so leaving it out is visible in the output.
\newcommand{\watermark}[1]{\gdef\grothWatermark{#1}}

\gdef\grothFolder{?}\gdef\grothBatch{}\gdef\grothFirst{?}\gdef\grothLast{?}%
\gdef\grothFoldertitle{}\gdef\grothShelfmark{}\gdef\grothDating{}\gdef\grothWatermark{}

% --- The résumé -------------------------------------------------------------
\newenvironment{resume}
  {\small\setlength{\parskip}{0.45em}}
  {\par\medskip\hrule\medskip}

% --- Keywords ---------------------------------------------------------------
% In English, closing the modernised reading's résumé: the modern vocabulary
% under which to search for what the leaves construct. The manifest extracts
% them to tag the volume — this line is the single source of the tags.
\newcommand{\keywords}[1]{\par\smallskip\noindent
  {\footnotesize\textsc{Keywords} --- \textit{#1}}\par}

% --- The critical apparatus -------------------------------------------------

% The Gallica view number, in the margin. This is the anchor that lets a
% disagreement over a formula be settled by looking at *one* image rather than
% a volume of 750. The site uses it to turn the facsimile as the transcription
% is scrolled. A view is one scanned image: usually one side of a leaf.
\newcommand{\page}[1]{%
  \marginnote{\footnotesize\color{gray!70}\textsf{v.\,#1}}%
  \label{page:#1}\ignorespaces}

% The BnF's pencilled foliation, where the leaf carries it and a pass can read
% it: « 198r ». This is what the literature cites — « ff. 198r–208v » — and
% the one line that turns a citation into a view. Recorded, never inferred:
% a folio a pass cannot read is not written.
\newcommand{\folio}[1]{%
  \marginnote{\footnotesize\color{gray!50}\textsf{f.\,#1}}\ignorespaces}

% The modernised reading's coarser counterpart to \page{}: which views a
% section is drawn from.
\newcommand{\pagerange}[2]{%
  \marginnote{\footnotesize\color{gray!70}\textsf{v.\,#1--#2}}%
  \ignorespaces}

% Illegible: never guessed.
\newcommand{\ill}{\textcolor{red!60!black}{[\,\ldots\,]}}

% A doubtful reading: the word is given, and flagged as uncertain.
\newcommand{\uncertain}[1]{\uline{#1}}

% An editorial addition — a missing letter, an implied word.
\newcommand{\add}[1]{\textcolor{blue!50!black}{[#1]}}

% Struck out by the writer. What was crossed out is part of the document.
\newcommand{\struck}[1]{\sout{#1}}

% The transcriber's note, on the state of the page or the mathematics.
\newcommand{\note}[1]{\par\smallskip{\small\color{gray!60!black}\textit{[#1]}}\par\smallskip}

% A marginal note of hers, distinct from the body of the text.
\newcommand{\marginal}[1]{\par\smallskip\noindent\hspace{1em}%
  {\small\color{gray!60!black}#1}\par\smallskip}

% --- Title ------------------------------------------------------------------
% The title block is French, like the documents themselves.

\AddToHook{shipout/background}{%
  \ifx\grothWatermark\empty\else
    \begin{tikzpicture}[remember picture, overlay]
      \node[rotate=52, scale=3.4, align=center, text=gray!18,
            font=\sffamily\bfseries]
        at (current page.center) {\grothWatermark};
    \end{tikzpicture}%
  \fi}

\AtBeginDocument{%
  \begin{center}
    {\large\bfseries Manuscrits de mathématiciens (Gallica) ---
      \ifx\grothShelfmark\empty\grothFolder\else\grothShelfmark\fi}\\[2pt]
    {\itshape\grothFoldertitle}\\[4pt]
    {\small vues \grothFirst--\grothLast
      \ifx\grothBatch\empty\else\ (lot \grothBatch)\fi}\\[2pt]
    \ifx\grothDating\empty\else
      {\small\color{gray!60!black}Datation du catalogue : \grothDating}\\[2pt]
    \fi
    \ifx\grothWatermark\empty\else
      {\small\bfseries\color{red!45!gray}\grothWatermark}\\[2pt]
    \fi
    {\footnotesize\color{gray!60!black}Transcription établie d'après les images IIIF de Gallica.
     Source gallica.bnf.fr / Bibliothèque nationale de France.\\
     Les lectures incertaines sont soulignées, les illisibles notées [\,\ldots\,],
     les ajouts entre crochets ; \textsf{v.} numérote les vues de Gallica,
     \textsf{f.} la foliotation au crayon de la BnF.}
  \end{center}
  \vspace{1em}}

\endinput


\folder{fr-9115}
\shelfmark{Français 9115}
\batch{20}
\pages{381}{400}
\dating{1801-1900}
\watermark{Lecture automatique\\non vérifiée}
\foldertitle{Papiers de Sophie GERMAIN. — Recueil de dissertations et problèmes mathématiques et physiques.}

\begin{document}

\note{Numérotation des feuillets. En haut à droite de chaque feuillet écrit sur son recto, un nombre à l'encre, d'une plume large : 186 (vue 381), 187 (383), 188 (385), 189 (387), 190 (389), 191 (391), 192 (392), 193 (394), 194 (396), 195 (398), 196 (400), sans lacune ; aucun sur les revers, pas même sur la vue 399, revers écrit du feuillet 195. Cette série occupe la place de la foliotation de la Bibliothèque, mais elle est tracée à la plume, comme dans les lots précédents ; aucune foliotation au crayon n'a été lue, et aucun « f. » n'est donc porté. Au milieu de la tête, entre parenthèses et barrés, les numéros de page de sa main : (31) à (35) aux vues 381 à 389 ; puis, pour le mémoire qui commence à la vue 398, (1) barré, (2) et (3). Les vues 382, 384, 386, 388, 390, 393, 395 et 397 sont des revers blancs où l'écriture ou les figures de l'autre face transparaissent à l'envers ; elles ne sont pas transcrites.}

\page{381}
\note{En tête, au milieu, un nombre entre parenthèses, lu « (31) », le second chiffre douteux, traversé d'un trait horizontal qui semble le biffer ; en haut à droite « 186 », à l'encre. Le feuillet commence au milieu d'une phrase dont le début est sur un feuillet antérieur, hors de ce lot.}

est fixée et l'autre appuyée, des valeurs relatives au mouvement dela lame dont les
deux extremités sont fixées; desorte que les six cas calculés par Euler peuvent etre
réduits à quatre, réellement independans les uns des autres, savoir: le cas I où les
deux extremités sont libres, § 15 du mémoire d'Euler; le cas III où une extremité
est fixée, et l'autre libre, § 30; le cas IV où les deux extremités sont appuyées, § 35;
enfin le cas VI où l'une et l'autre extremités sont fixées, § 43. On voit en effet
que la valeur des coëfficiens et l'équation qui détermine la valeur d'$\omega$ dans
le cas II d'une extremité libre et l'autre appuiée § 23, sont d'accord avec les
formules relatives au cas des deux extremités libres § 15; car on a dans le cas I,
lorsque sin.$\omega$ est positif, $\alpha = 1$, $\beta = -e^{\omega}$ $\gamma = 1 + e^{\omega}$ $\delta = 1 - e^{\omega}$; et l'équation delaquelle
on tire la valeur de l'angle $\omega$ est $\text{cos}\,\omega = \frac{2}{e^{\omega} + e^{-\omega}}$ que j'ai déja montré etre
alors equivalante à $\text{tang}\,\frac{1}{2}\omega = \frac{1 - e^{\omega}}{1 + e^{\omega}}$. on a demême, pour le cas II, $\alpha = 1$, $\beta = -e^{2\omega}$
$\gamma = 1 + e^{2\omega}$. $\delta = 1 - e^{2\omega}$; et léquation qui détermine la valeur de l'angle $\omega$, est
$\text{tang}\,\omega = \frac{1 - e^{2\omega}}{1 + e^{2\omega}}$. il est clair que la seule différence entre les formules tirées du
cas I et celles qui appartiennent au cas II, est que l'on a $\frac{1}{2}\omega$ dans les premieres,
au lieu de $\omega$ dans les secondes; cequi doit etre en effet, puisque, dans le
cas II, il s'agit d'une lame entière, et que, dans le cas I, la lame dont une
extremité est libre et l'autre appuyée, est considerée comme lamoitiée d'une
autre lame dont les deux extremités sont libres.\note{Les chiffres romains sont lus tels qu'ils sont écrits : « dans le cas II, il s'agit d'une lame entière » et « dans le cas I, la lame dont une extremité est libre et l'autre appuyée ». Les lignes précédentes du même feuillet appellent au contraire cas I celui des deux extrémités libres (§ 15) et cas II celui d'une extrémité libre et l'autre appuyée (§ 23) ; les deux chiffres paraissent intervertis dans cette phrase. Ils sont gardés.}

On devroit trouver les mêmes rapports entre le cas V d'une extremité fixée
et l'autre appuiée, § 39, et le cas VI des deux extremités fixées, § 43; cependant
on a dans le cas VI, lorsque sin.$\omega$ est positif, $\alpha = 1$. $\beta = -e^{\omega}$ $\gamma = -(1 + e^{\omega})$, $\delta = -(1 - e^{\omega})$
et $\text{cos}\,\omega = \frac{2}{e^{\omega} + e^{-\omega}}$ ou $\text{tang}\,\frac{1}{2}\omega = \frac{1 - e^{\omega}}{1 + e^{\omega}}$, tandis que, quoiqu'on ait àlaverité,
pour le cas VI, $\text{tang}\,\omega = \frac{1 - e^{2\omega}}{1 + e^{2\omega}}$, les coëfficiens, savoir $\alpha = 1$ $\beta = -1$, $\gamma = \pm\sqrt{2}\,(e^{2\omega} + e^{-2\omega})$
ne semblent avoir aucun rapport avec ceux du cas precedant; et il résulte
\note{Le premier « cas VI » de ce paragraphe (« on a dans le cas VI, lorsque sin.$\omega$ est positif ») est écrit ainsi, là où le sens attend le cas V ; les deux sont lus « VI ». Le radical de la dernière formule est suivi d'un signe lu « 2 » avec doute. Le feuillet s'arrête sur « il résulte » ; la suite est au feuillet suivant (vue 383). Le verso (vue 382) est blanc.}

\page{383}
\note{En tête, au milieu, « (32) » entre parenthèses, barré d'un trait horizontal ; en haut à droite « 187 », à l'encre. La première ligne continue la phrase qui clôt la vue 381 (« et il résulte »).}

delà que l'on a $y = \ldots\, e^{x\omega} - e^{-x\omega} \mp \sqrt{2}\,(e^{2\omega} + e^{-2\omega})\,\text{sin}\,x\omega$, aulieu qu'il semble que
l'on devroit avoir $y = \ldots\, e^{x\omega} - e^{(2-x)\omega} - (1 + e^{2\omega})\,\text{sin}\,x\omega - (1 - e^{2\omega})\,\text{cos}\,x\omega$.

Pour éclaircir cette difficulté, il faut observer que, dans l'analyse du cas V,
Euler suppose que l'extremité, pour laquelle $\underline{s}$ ou $x = 0$, est celle qui est appuyée,
tandis que, dans l'hypothese qui se rapporte au cas VI, ce ne peut être que pour
une des extremités dela lame entière et parconséquent pour l'extremité fixée de
la lame de longueur $\frac{1}{2}a$ que $x$ peut être supposé égal a zéro. et comme
il est fort indifferent d'adopter l'une ou l'autre de ces suppositions, on peut
choisir celle qui convient ici. Pour la transporter dans l'expression d'$y$, il
faut mettre \uncertain{1}$\,- x$ au lieu \add{de $x$} \struck{on} trouve ainsi
\[ y = \ldots\, e^{(1-x)\omega} - e^{(x-1)\omega} \mp \sqrt{2}\,(e^{2\omega} + e^{-2\omega})\,\text{sin}\,(1-x)\omega \]
\[ = \ldots\, e^{-\omega}\{e^{x\omega} - e^{(2-x)\omega} \pm \sqrt{2}\,(e^{2\omega} + e^{-2\omega})(\text{sin}\,\omega\,\text{cos.}\,x\omega - \text{sin}\,x\omega\,\text{cos.}\,\omega)\,e^{\omega}\} \]
\[ = \ldots\, e^{x\omega} - e^{(2-x)\omega} - (1 + e^{2\omega})\,\text{sin}\,x\omega - (1 - e^{2\omega})\,\text{cos}\,x\omega \]
en mettant pour sin.$\omega$ et cos.$\omega$ leurs valeurs $\pm\frac{e^{\omega} - e^{-\omega}}{\sqrt{2}\,(e^{2\omega} + e^{-2\omega})}$ $\pm\frac{e^{\omega} - e^{-\omega}}{\sqrt{2}\,(e^{2\omega} + e^{-2\omega})}$ (M. d'Euler § 39),
on parvient àla derniere forme de cette valeur qui est conforme à celle que
le rapport entre les cas V et VI fesait presumer.\note{Le chiffre qui précède « $-x$ » est couvert d'une tache d'encre, un chiffre écrit sur un autre ; « 1 » est proposé par le sens. « de $x$ » est écrit au-dessus de la ligne, et « on » de « on trouve » paraît recouvert par un signe « : ». Les deux fractions données pour sin.$\omega$ et cos.$\omega$ ont, sur le feuillet, le même numérateur $e^{\omega} - e^{-\omega}$ ; il est gardé tel quel. Après « fesait », quelques lettres très pâles, apparemment effacées, ne se lisent pas.}

On peut encore s'assurer que l'on est autorisé à considerer les lames dont
les deux extremités sont libres ou fixées, comme separés en deux autres moitiées
moins longues, dont une des extremités est appuiée dans le même point qui
appartient en même tems aumilieu des premieres lames, en observant que les
conditions indiquées par Euler, § \uncertain{13}, pour que l'extremité d'une lame
vibrante soit simplement appuyée, sont $y = 0$ et $\left(\frac{ddy}{ds^2}\right) = 0$,
et que lorsque sin.$\omega$ est positif, qui est le cas d'un nombre impair de
nœuds, la supposition $x = \frac{1}{2}$ donne, pour les lames dont il s'agit,
$y = (1 + e^{\omega})\,\text{sin}\,\frac{1}{2}\omega + (1 - e^{\omega})\,\text{cos}\,\frac{1}{2}\omega = 0$; d'ou il résulte $\left(\frac{ddy}{ds^2}\right) = -(1 + e^{\omega})\,\text{sin}\,\frac{1}{2}\omega - (1 - e^{\omega})\,\text{cos}\,\frac{1}{2}\omega = 0$.
\note{« une des extremités » : « des » et le début d'« extremités » sont récrits sur une première leçon, à l'encre plus noire. Le second chiffre du paragraphe d'Euler, « § 13 », a la queue descendante qui fait hésiter entre 3 et 9. La fin de la dernière formule, « $e^{\omega})$ cos », est empâtée d'encre et lue avec doute ; « $x$ » dans les exposants est la lettre écrite dans « $s$ ou $x = 0$ ». La ligne est la dernière du feuillet ; la suite est à la vue 385. Le verso (vue 384) est blanc.}

\page{385}
\note{En tête, au milieu, « (33) » entre parenthèses, barré d'un trait horizontal ; en haut à droite « 188 », à l'encre. Le feuillet s'ouvre sur un alinéa nouveau.}

On seroit porté par l'analogie à croire que, dans le cas des deux extremités libres,
lorsqu'il n'y a pas de nœud au milieu dela courbe, les deux moitiés dela lame
primitive pourroient etre également considerés comme des lames entieres soumises
aux mêmes conditions: mais les valeurs des ordonnés ne s'accomodent point à cette
hypothese; et on en voit la raison dans l'observation que j'ai faites, § \uncertain{14}, delaquelle
il resulte que la moitié dela partie vibrante comprise entre les deux nœuds les
plus voisins du milieu de la lame primitive, ne peut être égale aux parties
comprises entre les nœuds les plus voisins des extremités de cette lame et les mêmes
extremités

25. Les vibrations qui font l'objet du mémoire d'Euler et des remarques
précedentes, sont celles que M\textsuperscript{r} Chladni nomme vibrations transversales: mais les
vibrations des plaques elastiques libres, dans le cas où il ya des lignes nodales
suivant differentes directions, vibrations dont la théorie semble devoir être beaucoup
plus compliquée, presentent cependant une singularité remarquable; c'est que, tandis
que pour les vibrations \uncertain{transversales} dela lame elastique on ne peut, suivant cequ'on
vient devoir, \uncertain{supposer} que la lame entiere dont les extremités sont libres, soit
formé dela reunion de deux lames moitiées moins longues soumises aux mêmes
conditions, il ya aucontraire plusieurs cas où les plaques carrées élastiques vibrantes
libres de tous côtés peuvent êtres considerées comme composées de quatre autres
lames soumises aux mêmes conditions (voyez le traité d'Acoustique de M\textsuperscript{r} Chladni
§ \uncertain{110}). en comparant les différentes figures qui présentent entre elles le rapport
dont il s'agit, on voit que les figures composées\textsuperscript{+} sont celles dans lesquelles
il n'y a pas de lignes nodales qui coupent diamétralement la surface en
deux parties égals; et on sent aisement que cela doit être en effet,
puisque, pour qu'une partie \uncertain{contigue} puisse équivaloir à une extremité ou
\marginal{+ Les seules figures composées dont il soit question ici sont celles qui sont décomposables en quatre fig. simples. et il est clair que la même chose n'est pas vraie pour celles qui contiennent un nombre impair de fig. simples car alors il ya des lignes nodales diametrales}
\note{Le numéro 25 est écrit en tête de l'alinéa, dans la ligne ; le 5 a la forme d'un S. Le numéro du paragraphe renvoyé, « § 14 », est lu avec doute (le 4 ressemble à un h). « transversales » est écrit d'une manière où l'on croit lire « transvertales » ; « supposer » est récrit sur un premier mot, les lettres du milieu repassées. La note marginale est de sa main, d'une plume plus fine, dans la marge gauche en face des six dernières lignes ; elle est appelée par une croix après « composées » et commence elle-même par une croix double ; sa dernière ligne court au pied du feuillet sur toute sa largeur et s'arrête sur « diametrales » au bord du papier. Le texte principal s'arrête sur « à une extremité ou » ; la suite est à la vue 387. Le verso (vue 386) est blanc.}

\page{387}
\note{En tête, au milieu, « (34) » entre parenthèses, barré d'un trait horizontal ; en haut à droite « 189 », à l'encre. La première ligne continue la phrase qui clôt la vue 385 (« à une extremité ou »).}

coté libre, il faut qu'elle soit également libre et parconséquent en mouvement;
par exemple la figure \uncertain{63} de l'ouvrage de M\textsuperscript{r} Chladni qui est la plus simple
parmi celles qui presentent des lignes nodales en differens sens, répetée quatre
fois, forme, suivant la remarque de l'auteur, la fig. 68 a, figure àlaquelle
j'ai ajouté des lignes ponctuées pour marquer les limites des figures composantes.
on peut encore s'assurer, au moyen de lignes également ponctuées, que la même
figure \uncertain{63}, répeté neuf fois, donne la fig. 75, et la fig 82 lorsqu'elle est
prise seize fois. Ainsi, dans certains cas au moins, si on suppose une surface
vibrante élastique carré, dont la superficie soit exprimé par $aamm$, que l'on
divise ce carré en $mm$, autres carrés de grandeur $aa$, et que l'on applique
un stilet au point convenable, et sous les conditions propres à faire naitre dans
un des quatre carrés d'étendue $aa$ placés aux angles de la surface vibrante,
une certaine figure nodale; cette même figure sera répeté dans les $mm-1$,
autres carrés qui composent la surface entière; d'ou il résulte que le mouvement
de ces carrés élementaires est semblable dans chacun d'eux; et parconséquent/
que, si on suppose deux plans perpendiculaires entr'eux, et l'un et l'autre
perpendiculaires au plan horizontal, qui est en même tems celui de la surface
elastique en repos, les courbes formés par les intersections de ces plans supposés
encore perpendiculaires aux côtés de cette même surface elastique carrée, couperont
le plan horizontal dans des points dont les coordonnés seront $\alpha$, $\alpha + a$,
$\alpha + 2a$ \&c. et $\alpha'$, $\alpha' + a$, $2\alpha' + 2a$ \&c., $\alpha$ et $\alpha'$ étant dependans dela
nature de la figure nodale formé dans le carré de superficie $aa$, et deplus
que ces courbes seront symetriques autour des points dont il s'agit.

Cette propriété \add{que} les plaques carrées élastiques vibrantes libres de tous coté ont de
pouvoir etre considerées comme décomposées en un certain nombre de plaques
\note{Les numéros des figures renvoient à l'ouvrage de Chladni ; le premier, lu « 63 » deux fois, a un second chiffre qui peut être un 3 ou un 9 à queue courte. Les figures qu'elle dessine elle-même sont aux vues 391 et 396. « $mm-1$ » est suivi d'une virgule sur le feuillet. Dans la suite des coordonnées, le troisième terme de la seconde est écrit « $2\alpha' + 2a$ », là où la première suite fait attendre $\alpha' + 2a$ ; il est gardé. Dans la dernière phrase, « que » est écrit au-dessus de la ligne, « les » est récrit à l'encre plus noire sur un premier mot (apparemment « des »), et « ont » sur un autre (apparemment « est ») : la phrase a d'abord été « Cette propriété des plaques … est de pouvoir… ». Le feuillet s'arrête sur « de plaques » ; la suite est à la vue 389. Le verso (vue 388) est blanc.}

\page{389}
\note{En tête, au milieu, « (35) » entre parenthèses, barré d'un trait horizontal ; en haut à droite « 190 », à l'encre. La première ligne continue la phrase qui clôt la vue 387 (« en un certain nombre de plaques »).}

de même forme soumises aux mêmes conditions, propriété dont M\textsuperscript{r} Chladni donne
plusieurs exemples pour le cas ou $m = 2$, conduit donc à conclure qu'il existe dans certains
cas de grandes analogies entre les vibrations des plaques elastiques carrées et celles des
surfaces tendues demême figure: mais on peut voir, par la comparaison du § 24 du
mémoire de M\textsuperscript{r} Biot, inséré dans le tome 4 des memoires de l'institut, que, dans cedernier
genre de surfaces, $\alpha$ et $\alpha'$ sont zéro; cequi doit etre en effet, puisqu'alors les lignes
de repos sont necessairement placées aux limites des surfaces composantes. on voit
encore que, pour les figures composées de la fig. \uncertain{63}, $\alpha = \alpha' = \frac{1}{2}a$

La théorie complette des vibrations dela plaque carrée, comprendroit \uncertain{quinze}
cas differens, si on avoit égard à l'etat des côtés qui pourroient êtres appuyés,
libres ou fixés: mais il arriveroit sans doute, comme pour les vibrations \uncertain{transversales}
dela lame elastique, que ces cas auroient entr'eux des relations qui permettroient
de n'en calculer que quelques uns; par exemple, le cas où une figure régulière
formée par les lignes nodales dans une plaque carrée dont tous les côtés sont libres,
présente deux lignes diametrales qui se coupent au centre de cette figure,
pourrait peut-etre être ramené à celui dequatre plaques quarrées réunies dont
chacune auroit deux de ses côtés libres et les deux autres appuyés.

Je vais encore ajouter une reflexion sur la valeur du mot \emph{appuyé} employé dans
l'exposition de la théorie des lames vibrantes: il paroit qu'il s'applique à un état
fixe et déterminé, de façon à etre entierement equivallent aux nœuds de vibrations.
on conçoit en effet que, si la vibration étoit plus gênée par un point d'appui
artificiel qu'elle ne l'est par les autres nœuds, le mouvement ne pourroit pas
se transmettre àtravers cepoint; cequi doit pourtant avoir lieu, puisque l'on fait
abstraction dela largeur dela lame: mais lorsqu'il s'agit des plaques vibrantes,
je ne sais si on peut affirmer la même chose; car alors le mouvement pourroit
etre transmis par les points adjacens; et l'irregularité de certaines figures observées par
M\textsuperscript{r} Chladni me porte à croire qu'il ya des cas où le mouvement est plus gêné dans une premiere
moitié dela plaque que dans la
\note{Le chiffre du tome des mémoires de l'Institut et celui du § 24 ont la même forme, un 4 tracé comme un h. « deux », dans « chacune auroit deux de ses côtés », est écrit sur un premier mot et empâté. « appuyé » est souligné. « quinze » est écrit « quize » en apparence, le n à peine formé. Les derniers mots, depuis « dans une premiere », sont serrés en deux lignes d'une écriture plus petite contre le bord droit du feuillet, et s'arrêtent sur « dans la » ; la suite est à la vue 391 ou plus loin, voir la note de cette vue. Le verso (vue 390) est blanc.}

\page{391}
\note{Feuillet de figures, sans texte suivi ; en haut à droite « 191 », à l'encre ; pas de numéro entre parenthèses en tête. Il porte six figures à la plume. (1) En haut, « fig \uncertain{1} » : une lame $EF$ droite et, par-dessus, sa courbe de vibration, qui passe au-dessus de l'axe en $m$, le coupe en $L$ et passe au-dessous en $n$ avant de revenir en $F$. (2) Au-dessous, « fig. \uncertain{1}$'$ » : la même lame $EF$ avec une courbe à sept arcs qui coupe l'axe en $L$, $L'$, $L''$, $L'''$, $L^{\text{IV}}$, $L^{\text{V}}$, $L^{\text{VI}}$, les arcs alternativement au-dessus et au-dessous. Le chiffre de ces deux légendes a le crochet initial du 1 du nombre « 191 » du même feuillet, et peut aussi se lire 7. (3) « fig. \uncertain{63} » : un carré partagé en quatre par deux lignes nodales qui se croisent au centre, le croisement noirci en étoile. (4) « fig. 68 a » : un carré partagé par deux verticales et deux horizontales en traits pleins, avec, en pointillé, la médiane verticale et la médiane horizontale ; les quatre points où les traits pleins se croisent sont noircis. C'est la figure 63 répétée quatre fois, les lignes ponctuées marquant « les limites des figures composantes » (vue 387). (5) « fig. 75 » : le même dessin sur un carré partagé en neuf, avec deux verticales et deux horizontales ponctuées et neuf croisements noircis. (6) « fig 82 » : le même encore, sur un carré partagé en seize, trois lignes ponctuées dans chaque sens. Les numéros 63, 68 a, 75 et 82 sont ceux que cite le texte de la vue 387 (« l'ouvrage de M\textsuperscript{r} Chladni ») ; le 3 de « 63 » a une queue qui le fait ressembler à un 9. La vue suivante (392) n'est pas le revers de ce feuillet mais un autre feuillet, plus petit, numéroté 192 ; aucune vue de ce lot ne montre le revers de celui-ci.}

\page{392}
\note{Un feuillet plus petit que les précédents, écrit sur sa moitié supérieure ; en haut à droite « 192 », à l'encre ; pas de numéro entre parenthèses en tête. Ce sont des calculs sans phrase d'introduction : rien n'y enchaîne sur la fin de la vue 389 (« que dans la »).}

\[ e^{x\omega} \pm e^{\omega - x\omega} + (1 \mp e^{\omega})\,\text{sin}\,x\omega + (1 \pm e^{\omega})\,\text{cos}\,x\omega = 0 \]
soit \[ e^{\frac{1}{2}x\omega} \pm e^{\omega - \frac{1}{2}x\omega} + (1 \mp e^{\omega})\,\text{sin}\,\tfrac{1}{2}x\omega + (1 \pm e^{\omega})\,\text{cos}\,\tfrac{1}{2}x\omega = 0 \]
\[ 2\,\text{cos}\,\tfrac{1}{2}x\omega\,(e^{\frac{1}{2}x\omega} \pm e^{\omega - \frac{1}{2}x\omega}) = e^{x\omega} \pm e^{\omega - x\omega} + 1 \pm e^{\omega} \]
\[ = (e^{x\omega} + 1)(1 \pm e^{\omega - x\omega}) \]
\[ 2\,\text{cos}\,\tfrac{1}{2}x\omega = (e^{\frac{1}{2}x\omega} + e^{-\frac{1}{2}x\omega}). \]
resultat absurde, on voit que l'on y est parvenu en retranchant
la seconde équation multipliée par $2\,\text{cos}\,\frac{1}{2}x\omega$ dela premiere
il faut donc que l'équation $e^{\frac{1}{2}x\omega} \pm$ \uncertain{\&c} elle même soit inadmissible
or cette
\note{Le feuillet s'arrête sur « or cette », au premier tiers de sa hauteur ; le reste est blanc. Le signe lu « \&c » après l'exponentielle de l'avant-dernière ligne est une boucle suivie d'un trait, lue avec doute. « resultat » commence par une minuscule. Dans la deuxième équation, le second exposant est lu $\omega - \frac{1}{2}x\omega$, comme dans la troisième.}

\page{394}
\note{Un petit feuillet du même format que celui de la vue 392, écrit sur sa moitié supérieure ; en haut à droite « 193 », à l'encre ; pas de numéro entre parenthèses en tête. Le cachet « Bibliothèque Impériale MSS. » est frappé à droite, contre les deux dernières lignes. Mêmes calculs que la vue 392, dans les mêmes notations.}

si l'ordonnée qui repond a $\frac{1}{2}x$ etoit laplus grande et que
la tangente a cepoint fut parrallele a l'axe c'est adire
si $\left(\frac{dy}{dx}\right)$ etoit egal a zéro il en resulteroit l'équation
\[ e^{\frac{1}{2}x\omega} \mp e^{(1-\frac{1}{2}x)\omega} + (1 \mp e^{\omega})\,\text{cos}\,\tfrac{1}{2}x\omega - (1 \pm e^{\omega})\,\text{sin}\,\tfrac{1}{2}x\omega = 0 \]
ou \[ (e^{\frac{x}{2}\omega} \mp e^{(1-\frac{1}{2}x)\omega})\,2\,\text{sin}\,\tfrac{1}{2}x\omega + (1 \mp e^{\omega})\,\text{sin}\,x\omega + (1 \pm e^{\omega})(\text{cos}\,x\omega - 1) = 0 \]
\[ 1 \pm e^{\omega} - (e^{\frac{1}{2}x\omega} \mp e^{(1-\frac{1}{2}x)\omega})\,2\,\text{sin}\,\tfrac{1}{2}x\omega = (1 \mp e^{\omega})\,\text{sin}\,x\omega + (1 \pm e^{\omega})\,\text{cos}\,x\omega \]
on a d'ailleurs
\[ e^{x\omega} \pm e^{\omega - x\omega} + (1 \mp e^{\omega})\,\text{sin}\,x\omega + (1 \pm e^{\omega})\,\text{cos}\,x\omega = 0 \]
ainsi on auroit
\[ 1 + e^{x\omega} \pm e^{\omega} \pm e^{\omega - x\omega} = 2\,(e^{\frac{1}{2}x\omega} \mp e^{(1-\frac{1}{2}x)\omega})\,\text{sin}\,\tfrac{1}{2}x\omega \]
\[ (1 + e^{x\omega})(1 \pm e^{(1-x)\omega}) = 2\,(e^{\frac{1}{2}x\omega} \mp e^{(1-\frac{1}{2}x)\omega})\,\text{sin}\,\tfrac{1}{2}x\omega \]
\note{La première ligne commence par une minuscule, « si ». Dans la deuxième équation, le premier exposant est écrit $\frac{x}{2}\omega$ et non $\frac{1}{2}x\omega$. Le second exposant de la dernière ligne est empâté. Le feuillet s'arrête là, aux deux tiers de sa hauteur ; le reste est blanc. Le revers (vue 395) est blanc : on y voit par transparence, à l'envers, les calculs de ce feuillet et le cachet.}

\page{396}
\note{Second feuillet de figures, sans texte ; en haut à droite « 194 », à l'encre, le dernier chiffre contre le bord du papier ; le cachet « Bibliothèque Impériale MSS. » au milieu du feuillet. Six figures de plaques carrées à la plume, chacune tracée sur un quadrillage en pointillé, avec ses lignes nodales en traits pleins de plusieurs épaisseurs. (1) « fig. 67 c » : un carré $a\,c\,b\,d$ (a et c en haut, b et d en bas), la diagonale $a\,d$ en trait fort, et une courbe fermée arrondie qui traverse le carré dans l'autre sens ; des petites lettres $e'$ et $i$ le long du quadrillage. (2) « fig 72 a » : le même carré $a\,c\,b\,d$, les deux diagonales, et une courbe fermée arrondie plus grande ; mêmes petites lettres $e$, $i$. (3) « fig 78 » : un carré $a$, $m$, $c$ en haut, $n$ et $K$ au milieu des côtés, $v$ sur le côté droit, $b$, $h$, $a$, $d$ en bas, $g$ au centre ; une diagonale et des courbes nodales en festons, plusieurs d'entre elles doublées d'un trait hachuré de petites croix. (4) « fig. 85 » : un carré $a\,b\,c\,d$ ($a$ et $b$ en haut, $c$ et $d$ en bas), $n$ et $K$ aux milieux des côtés, les deux diagonales, deux boucles en haut et, dans la moitié inférieure, des courbes ondulées doublées de croix. (5) Sans numéro : un carré aux deux diagonales, la demi-diagonale supérieure gauche en trait très fort, le centre $C$ noirci ; $A$ et $F$ à gauche, $D$ sur le côté droit, $B$ et $E$ en bas ; de $C$ partent un trait vers $B$ et un trait vers $D$, et le petit carré $C\,D\,E\,B$ est fermé en tirets avec sa diagonale $C\,E$ en tirets. (6) Sans numéro : un rectangle plus large que haut, $b$ et $d$ en haut, $b$ et $d$ en bas, sur un quadrillage pointillé, avec deux diagonales croisées en $M$ et une courbe en forme de cœur, les croisements marqués $n$ ; mêmes petites lettres $i$, $c'$. Les numéros 67 c, 72 a, 78 et 85 ne sont cités par aucun texte de ce lot. Le revers (vue 397) est blanc : on y voit par transparence, à l'envers, les figures de ce feuillet.}

\page{398}
\note{Un nouveau morceau commence ici. En tête, au milieu, « (1) » entre parenthèses, barré d'un trait oblique, la parenthèse fermante mal formée ; en haut à droite « 195 », à l'encre. Le cachet « Bibliothèque Royale » (couronne, lettre M) est frappé dans la marge gauche, en face du titre. C'est un brouillon très corrigé : lignes entières barrées d'un trait épais, la leçon nouvelle écrite dans l'interligne au-dessus ou au-dessous, des renvois par croix.}

Mémoire sur \struck{cette} la question proposée par la premiere classe de
l'institut
« Donner la théorie mathématique des vibrations des
« surfaces élastiques et la comparer à l'experience ».

\begin{quote}
Sed longe maximum progressibus scientiarum, et novis pensis ac provinciis in isdem suscipiendis,
obstaculum deprehenditur, in desperatione hominum et suppositione impossibilitis. \emph{Baco, novum organum}, lib. 1
n.\textsuperscript{o} XCII. tome 1 p. 298. edit in f\textsuperscript{o} \uncertain{lond}\ill{}
\end{quote}
\note{La citation latine est écrite d'une plume plus fine, serrée entre le titre et la première ligne du texte, et s'achève à droite en deux lignes au-dessus de « d'abord la question » ; la référence est soulignée. « isdem » et « impossibilitis » sont lus tels qu'ils sont écrits. Le dernier mot de la référence commence par « lond » ; la fin ne se lit pas.}

Si on vouloit \add{examiner} \struck{embrasser} d'abord \struck{cette} laquestion dans toute sa généralité,
on tomberoit dans des calculs extrèmement compliqués; mais en se
bornant à considerer le mouvement des surfaces planes dans le cas
où il est regulier et comparable à celui d'un pendule simple, on parvient
aisément à une équation \add{differentielle dont les intégrales particulieres sont facilement applicables}\struck{fort simple}. J'ose esperer que, malgré les
simplifications que, \struck{conformement} \add{d'après} \add{a cette remarque,} j'ai cru pouvoir introduire dans mes formules,
la classe trouvera la théorie que \struck{j'ai} je \struck{l'honneur de} soumets a son
jugement, aussi générale que le demande l'état de la question,
car, d'une part, les experiences de M\textsuperscript{r} Chladni sont toutes comprises
dans le cas des mouvemens réguliers, les seuls en effet qui appartiennent
à l'acoustique\textsuperscript{+}; \add{+ \struck{en tant qu'elle s'occupe des lois relatives au son}} \struck{de l'autre, cet habile physicien \uncertain{n'a fait qu'un}
petit nombre d'essais sur le son des surfaces courbes; et d'ailleurs
les géomètres sentiront aisément que} \add{et de l'autre} les simplifications auxquelles
\struck{donne lieu} la supposition de $\underline{z}$ très petit \add{+ donne lieu} \struck{pourroient} \add{peuvent aisément} se transporter
aux \add{petit nombre de mouvemens que cet habile physicien a observé dans} \struck{surfaces d'une courbure quelconque en substituant \ill{}}
les surfaces courbes car, comme je le dirai dans la suite, ces
\struck{différentielles \ill{} les différentielles}
mouvemens se trouvent compris dans une équation fort
\struck{de la courbe \ill{}}
simple qui convient à plusieurs genres de corps sonores et dans laquelle
\struck{desorte que les valeurs de} $\underline{z}$ exprime\struck{roient \ill{}}, comme dans le
cas du plan, la quantité dont chaque point des mêmes
surfaces s'élev\struck{oit} pendant leur mouvement audessus de \struck{leur} sa
situation naturelle.
\marginal{puisqu'eux seuls presentent des lois \struck{qu'elle \ill{} appréciables} déterminables;}
\note{L'ordre des corrections est celui que la page suggère ; il est offert, non assuré. (1) Une première addition interlinéaire, précédée d'une croix, « en tant qu'elle s'occupe des lois relatives au son », est barrée ; la note marginale, sans signe de renvoi, écrite en face de ces lignes, paraît la remplacer : « puisqu'eux seuls presentent des lois déterminables », « déterminables » étant écrit sous « qu'elle … appréciables » barrés, un pâté entre les deux mots. (2) La phrase barrée « de l'autre, cet habile physicien … » tient une ligne et demie ; « et de l'autre » est récrit au-dessus de la fin de la ligne suivante. (3) « donne lieu », barré en tête de ligne, est récrit dans l'interligne après « très petit », précédé d'une croix, suivi de « peuvent aisément » qui remplace « pourroient », barré. (4) « petit nombre de mouvemens que cet habile physicien a observé dans » est écrit au-dessus de « surfaces d'une courbure quelconque en substituant … », barré d'un trait épais. (5) Trois lignes barrées d'un trait épais alternent avec les lignes gardées, qui ont été écrites entre elles ; on n'y lit que « différentielles » au début et à la fin de la première, et « de la courbe » au début de la deuxième ; la troisième commence par « desorte que les valeurs de ». La lettre soulignée, deux fois, est lue $z$ ; elle ressemble à un 2. « s'élevoit » a sa terminaison barrée. Au pied du feuillet, des calculs à l'envers transparaissent : ce sont les formules du bas de la page (2), écrite au revers de ce feuillet et photographiée à la vue 399, où le texte continue.}

\page{399}
\note{Revers du feuillet de la vue 398, écrit : en tête, au milieu, « (2) », non barré ; pas de nombre à l'encre en haut à droite, le revers n'en portant pas. Le texte fait suite au préambule de la page (1), vue 398, et ouvre le premier paragraphe du mémoire. Les calculs du bas de cette page sont ceux qu'on voit transparaître, à l'envers, au pied de la vue 398.}

§ 1. Recherche de l'équation différentielle dela surface élastique vibrante

N.\textsuperscript{o} 1. Pour trouver cette équation, il faut d'abord déterminer quelle
doit être la somme des momens des forces d'elasticité qui agissent
dans toute l'étendue dela surface. Pour cela, je prends
quatre points sur cette surface; j'imagine deux plans, l'un
passant par le premier, le second et le quatrième de ces
points, l'autre passant par le premier, le troisième et le
quatrième des mêmes points; je nomme $\underline{e}$ l'angle extérieur
formé par un de ces plans et le prolongement de l'autre;
je represente par $\underline{\beta}$ la petite ligne d'intersection des deux plans
et par $E$ la force qui agit suivant $\beta$ en s'opposant à
l'inflexion de la surface et en tendant par conséquent à
diminuer l'angle $\underline{e}$. \struck{d'où il suit} \add{Ainsi je trouve} que la somme des momens
est exprimée par $\int\!\!\int E\beta\,de$; et \struck{qu'} il ne s'agit plus, pour
avoir sa valeur, que de connoître celles de $de$, $\underline{\beta}$ et $\underline{E}$.

Je commence par la recherche dela valeur de $e$.
On sait qu'en représentant par $\underline{C}$ l'ordonnée verticale
d'un plan à son origine, par $A$ et $B$ les tangentes des
angles que les traces du plan dont il s'agit sur les plans
coordonnés verticaux font avec l'horizon, et par $C'$, $A'$ et $B'$
les quantités analogues relatives à un second plan, le cosinus
de l'angle que ces deux plans forment entr'eux, est
exprimé par $\frac{AA' + BB' + CC'}{\sqrt{A^2 + B^2 + C^2}\,.\,\sqrt{A'^2 + B'^2 + C'^2}}$. On a donc ici
\[ \text{cos}\,(\pi - e) = -\,\text{cos}\,e = \frac{AA' + B'B' + CC'}{\sqrt{A^2 + B^2 + C^2}\,.\,\sqrt{A'^2 + B'^2 + C'^2}} \]
d'où on tire
\note{Le titre du paragraphe est souligné d'un trait sur toute sa longueur ; son numéro, en tête, est écrit avec le même signe § que dans les pages précédentes. « d'où il suit » est barré et « Ainsi je trouve » écrit au-dessus ; « qu' », après « et », est barré d'un trait épais. La lettre de la « petite ligne d'intersection » est tracée comme le $\beta$ des coefficients d'Euler aux vues 381–383, une boucle montante sur la ligne ; elle est rendue $\beta$. Dans la seconde fraction, le deuxième terme du numérateur est écrit « $B'B'$ », au lieu de $BB'$ de la première ; il est gardé. Un pâté d'encre au début de la ligne « quatrième des mêmes points ». La page s'arrête sur « d'où on tire » ; la suite est à la vue 400.}

\page{400}
\note{En tête, au milieu, « (3) », non barré ; en haut à droite « 196 », à l'encre. La page fait suite à la page (2) (vue 399), qui s'arrête sur « d'où on tire ».}

\[ \text{sin}\,e^2 = \frac{(A^2 + B^2 + C^2)(A'^2 + B'^2 + C'^2) - (AA' + BB' + CC')^2}{(A^2 + B^2 + C^2)(A'^2 + B'^2 + C'^2)} \]
\[ = \frac{(AB' - A'B)^2 + (AC' - A'C)^2 + (BC' - B'C)^2}{(A^2 + B^2 + C^2)(A'^2 + B'^2 + C'^2)} \]
On a encore par les formules connues
\[ A = z'(y'' - y''') + z''(y''' - y') + z'''(y' - y'') \]
\[ B = x'(z'' - z''') + x''(z''' - z') + x'''(z' - z'') \]
\[ C = y'(x'' - x''') + y''(x''' - x') + y'''(x' - x'') \]
équations dans lesquelles $x', y', z'$; $x'', y'', z''$; $x''', y''', z'''$, sont les coordonnés
des trois points par lesquels le plan est supposé passer. Pour déterminer
les valeurs de ces coordonnés dans le cas présent, rien n'empêche de
supposer qu'en passant d'un point pris sur la surface à un
autre point infiniment voisin du premier, une des coordonné soit
regardée comme constante, tandis que les deux autres seront supposés
variables. Ainsi on pourra prendre $x, y, z$; pour les coordonnés du
premier point $\ldots\ldots\ldots\;\; x,\ y + dy,\ z + dz$; pour les
coordonnés du second point $\ldots\;\; y,\ x + dx,\ z + dz$; pour les
coordonnés du troisième point, enfin $x + dx,\ y + dy,\ z + 2dz + d^2z$
pour les coordonnés du quatrième point

J'ai supposé que le premier plan passoit par le premier, le second
et le quatrième point. On a donc pour ce plan
\[ x' = x,\quad y' = y,\quad z' = z;\quad x'' = x,\quad y'' = y + dy,\quad z'' = z + dz;\quad x''' = x + dx, \]
\[ y''' = y + dy,\quad z''' = z + 2dz + d^2z \]
\[ A = z'(y'' - y''') + z''(y''' - y') + z'''(y' - y'') = z\,.\,0 + (z + dz)\,dy - (z + 2dz + d^2z)\,dy \]
\[ = -(dz + d^2z)\,dy \]
\[ B = x'(z'' - z''') + x''(z''' - z') + x'''(z' - z'') = x(-dz - d^2z) + x(2dz + d^2z) - (x + dx)\,dz = -dz\,dx \]
\[ C = y'(x'' - x''') + y''(x''' - x') + y'''(x' - x'') = -y\,dx + (y + dy)\,dx + (y + dy) \times 0 = dx\,dy \]
\note{Le dernier terme de la première formule est lu $(BC' - B'C)^2$ ; les accents y sont mal placés, tous deux au-dessus du C, et la lecture est offerte avec doute. Les lignes de points après « premier point » et « second point » sont d'elle ; elles alignent les coordonnées en colonne. $d^2z$ est écrit $d$ puis $z$ surmonté d'un 2. Dans la ligne de $A$, le $y'$ du dernier facteur est récrit et empâté ; dans celle de $B$, le premier $x$ après le signe égal est empâté. La page est pleine ; la suite est à la vue 401, hors de ce lot.}

\end{document}
